% Suggested replacement prose only; not installed in paper_MFP/main.tex.
% Citation keys are supplied in references.bib in this study.
\section{Introduction}

How can one calculate the expected initial derivatives of a deep neural
network without deriving a separate system of equations for each observable?
Even at fixed depth and fixed derivative order, the calculation rapidly
involves forward activations, backward vectors, products of derivatives,
normalized contractions, and repeated uses of the same weight matrices.
Training derivatives add a further dependency: the direction of motion
changes with the parameters and with the predictions on the entire batch.
A useful calculus must retain these dependencies while turning the requested
quantity into an explicit computation whose dimension is independent of
network width.

The probabilistic foundation for such calculations is supplied by Tensor
Programs. Its Gaussian source-response rules account for matrix and transpose
reuse, including singular query covariances and causal scalar averages
\cite{yang2021tp3}. Smooth-program moment estimates supply the corresponding
expectation and finite-$L^p$ limits \cite{golikov2022nongaus}. The later
Tensor Programs IVb develops a broader language and an explicit
backpropagation calculus \cite{yang2023tp4b}. These results
already support feature-learning limits under maximal-update
parameterization \cite{yang2021tp4}. Higher-order neural derivative
hierarchies and symbolic Gaussian calculations also have substantial
precedents \cite{huang2020nth,chen2020lantk,dyer2019feynman,aitken2020nonlinear}.
The question addressed here is constructive: how to specify a derivative
query, compile it without changing the finite network it describes, and
return an inspectable Gaussian-expectation formula with a precise domain of
validity.

Mean-field peeling (MFP) provides this construction for a declared typed
language of finite Gaussian neural programs. The language contains
coordinate nonlinearities and their derivatives, sums and products,
normalized averages, named matrix and transpose calls, and represented
parameter updates. Its differentiation rules cover the specified input and
parameter directions, scaled gradients, admitted contractions, finite
updates, and finite moving-flow jets. Admissibility is determined by these
rules; it is not an assumption that an arbitrary derivative tensor can be
reduced successfully.

For an admitted scalar observable $O_n$ at width $n$, the construction
returns a finite Gaussian-expectation graph $\mathcal G_O$ such that
\[
  O_n \longrightarrow \operatorname{Eval}(\mathcal G_O)
  \quad\text{in }L^p,\qquad 1\leq p<\infty.
\]
The program, depth, sample count, derivative order, and number of updates
are fixed before $n\to\infty$. Initial root tuples are independent across
neurons, with fixed Gaussian laws that may have singular covariances;
initial scalar parameters are deterministic. The activation is smooth and
every derivative has polynomial growth. The theorem also supplies uniform
moments and joint empirical laws for the resulting vector computations.

There are two exactness statements. Compilation by the product, chain and
adjoint rules is an identity at every finite width. Gaussian evaluation is
an exact formula for the resulting width limit. For example, with
$W_{ij}\sim N(0,1/n)$ independently and $G\sim N(0,1)$,
\[
 \lim_{n\to\infty}\mathbb E
 \left[\frac1n\left\|W^{\mathsf T}\phi(W\mathbf1)\right\|^2\right]
 =\mathbb E[\phi(G)^2]+\mathbb E[\phi'(G)]^2.
\]
The second term records reuse of the original matrix. For identity
activation the finite-width expectation is $2+1/n$, while the graph
evaluates to $2$. The same distinction applies to more elaborate neural
queries.

The construction separates physical derivatives of the finite network from
formal Gaussian-source derivatives in the response corrections. Physical
derivatives propagate through computed scalars and moving directions.
Formal source derivatives hold already evaluated scalar expectations fixed
and retain distinct source slots even when covariance is singular. These
conventions determine the executable rules, rather than being left to
informal simplification.

The accompanying symbolic compiler implements this interface and returns
the scalar formula, its Gaussian integrands and covariance tables, and the
transformation trace. Generic Gaussian integrals may remain unevaluated;
polynomial specializations can reduce to exact arithmetic. Worked examples
progress from a reused transpose and represented updates to MLP gradient
observables and the initial jets of a final hidden feature kernel under
batch gradient flow, with symbolic input-geometry factors and labels.
The latter kernel is a specified feature kernel, not the full tangent
kernel. A separate initialization result reduces a restricted class to
activation-derivative moments by Wick and Stein identities; its literal
preactivation restriction is essential.

The contribution is this explicit derivative language, constructive
correctness argument, symbolic implementation, and collection of fully
specified calculations. It is not a new general probabilistic program
limit theorem. Nor does finiteness of the output graph imply efficient
compilation, elementary Gaussian integrals, convergence of the Taylor
series, or reconstruction of the network at positive training time.
These boundaries keep the calculus usable without attributing guarantees
to it that its finite-query theorem does not provide.
